% ======================================================================
%  UCL MSc Dissertation Template
%  Based on the UCL postgraduate project report template
%  (project_report_-_pg.tex), extended with the front matter and
%  chapter structure of a master's dissertation.
%
%  Compile with:  pdflatex main  ->  bibtex main (optional)
%                 -> pdflatex main x2
% ======================================================================

\documentclass[a4paper,12pt]{report}

% ---------- Packages -------------------------------------------------
\usepackage{setspace}            % line spacing
\usepackage{amssymb,graphicx,color}
\usepackage{amsfonts}
\usepackage{latexsym}
\usepackage{a4wide}
\usepackage{amsmath}
\usepackage[round]{natbib}       % author-year (Harvard) citations; load BEFORE hyperref
\usepackage[hidelinks]{hyperref} % clickable TOC / refs
\usepackage{booktabs}            % professional tables
\usepackage{caption}
\usepackage{subcaption}          % subfigures (modern replacement for subfigure)
%\usepackage[backend=bibtex]{biblatex} % alternative bibliography backend

\pagestyle{plain}

% ---------- Theorem environments -------------------------------------
\newtheorem{theorem}{THEOREM}
\newtheorem{lemma}[theorem]{LEMMA}
\newtheorem{corollary}[theorem]{COROLLARY}
\newtheorem{proposition}[theorem]{PROPOSITION}
\newtheorem{remark}[theorem]{REMARK}
\newtheorem{definition}[theorem]{DEFINITION}
\newtheorem{fact}[theorem]{FACT}
\newtheorem{problem}[theorem]{PROBLEM}
\newtheorem{exercise}[theorem]{EXERCISE}

\def \set#1{\{#1\} }

\newenvironment{proof}{%
PROOF:
\begin{quotation}}{%
$\Box$ \end{quotation}}

% ---------- Handy macros ----------------------------------------------
\newcommand{\nats}{\mbox{\( \mathbb N \)}}
\newcommand{\rats}{\mbox{\(\mathbb Q\)}}
\newcommand{\reals}{\mbox{\(\mathbb R\)}}
\newcommand{\ints}{\mbox{\(\mathbb Z\)}}
\newcommand{\Gle}{\ensuremath{\mathcal{G}_{\leq t}}} % graph of edges with timestamp <= t

% ---------- Dissertation metadata (EDIT THESE) ------------------------
\newcommand{\thesistitle}{Inductive Temporal Theme Velocity Discovery on a Knowledge Graph}
\newcommand{\thesisauthor}{Denis Maruev}
\newcommand{\thesisdegree}{MSc Computational Statistics and Machine Learning}
\newcommand{\thesissupervisor}{Prof.\ Philip Treleaven}
\newcommand{\thesisadvisor}{Dr.\ Olga Kolchyna}      % leave empty if none
\newcommand{\submissiondate}{Day Month Year}

% ======================================================================
\begin{document}
\onehalfspacing

% ---------------------------------------------------------------------
% TITLE PAGE
% ---------------------------------------------------------------------
\begin{titlepage}
\begin{center}
% Upload ucl_logo.png to your Overleaf project, then uncomment the next line:
%\includegraphics[scale=.5]{ucl_logo.png}\\[1.5cm]

{\Huge \thesistitle}\\[0.4cm]
% Uncomment the next line if you have a subtitle:
%{\large Optional Subtitle}\\[0.4cm]
\vspace{2cm}

{\Large \thesisauthor}\\[1.5cm]

A dissertation submitted in partial fulfilment\\
of the requirements for the degree of\\[0.3cm]
{\bf \thesisdegree}\\
of\\
{\bf University College London}\\[2cm]

Department of Computer Science\\
University College London\\[1cm]

Academic Supervisor: \thesissupervisor\\
Technical Advisor: \thesisadvisor\\[1cm]

Submission date: \submissiondate
\end{center}

\vfill
\footnotesize
{\bf Disclaimer:}
This report is submitted as part requirement for the \thesisdegree\ at UCL.
It is substantially the result of my own work except where explicitly
indicated in the text.
\emph{Either:} The report may be freely copied and distributed provided the
source is explicitly acknowledged.
\emph{Or:} The report will be distributed to the internal and external
examiners, but thereafter may not be copied or distributed except with
permission from the author.
\end{titlepage}

% ---------------------------------------------------------------------
% FRONT MATTER
% ---------------------------------------------------------------------
\pagenumbering{roman}

\chapter*{Declaration}
\addcontentsline{toc}{chapter}{Declaration}
I, \thesisauthor, confirm that the work presented in this dissertation is my
own. Where information has been derived from other sources, I confirm that
this has been indicated in the dissertation.

\vspace{2cm}
\noindent\rule{6cm}{0.4pt}\\
\thesisauthor

\begin{abstract}
\addcontentsline{toc}{chapter}{Abstract}
% ===================== ABSTRACT (inlined) ============================
A knowledge graph represents entities and the relationships between them; a
time-indexed knowledge graph additionally records when each relationship forms,
making its dynamics measurable. This dissertation studies the dynamics of a
financial knowledge graph built from public news: how fast relationships between
companies, sectors, countries and themes form (\emph{velocity}), how that rate
changes (\emph{acceleration}), and whether that signal predicts what forms next.
The hypothesis is that the rate at which entities become connected leads which
connections form, treated as a graph-native question under strict point-in-time
integrity. The work comprises three experiments, which constitute its original
research and are discussed in Section~1.3.

\textbf{Experiment 1: Building the Time-Indexed Knowledge Graph.} The substrate
is constructed by running the open-source ICKG extractor (ICKG-v4.2, a LoRA
adapter on Qwen-2.5-14B-Instruct distilled from GPT-4) over the FNSPID full-text
news corpus, producing FinDKG-compatible quadruples (subject, relation, object,
time) over twelve entity and fifteen relation types, with themes realised as
CONCEPT and EVENT nodes and every edge timestamped by publication date. Two
corpus scales are built (201{,}600 and 604{,}800 time-stratified articles; 0\%
malformed output; $\sim$18.5 valid triples per article), and a syndication audit
shows that 10--24\% of corpus articles are verbatim duplicates, motivating a
deduplicated core of 21{,}216 entities and 212{,}294 edge-weeks that serves as
the canonical graph.

\textbf{Experiment 2: Velocity, Bursts, and Walk-Forward Theme Discovery.} This
experiment --- the dissertation's central contribution --- defines and
measures the rate and acceleration of new-edge
formation per node, relation and theme (trailing-window, past-only
Poisson-surprise normalisation), detects bursts with three independent detectors
(Kleinberg's automaton, the CUSUM control chart, and Bayesian online
changepoint detection), and clusters
co-accelerating entities on their recent co-occurrence subgraph into
\emph{themes}, tracked week-over-week into lifelines whose birth dates are the
dates the system would have flagged each theme live. Against a frozen
ten-event gold list the system detects 8/10 major 2020--2023 themes at a median
lead of $-1$~week --- anticipatable themes early (the COVID market crisis
$-3$~weeks, the vaccine race $-8$) and pure shock events at the news-flow floor
($+1$) --- and a 100{,}000-draw placebo control shows that random event weeks
would score $\sim$2/10 ($p \approx 10^{-4}$).

\textbf{Experiment 3: Temporal Link Prediction under Point-in-Time Discipline.}
The third experiment maps what is and is not forecastable at the
individual-edge level. An RE-GCN-style temporal graph network is evaluated
against the static embeddings ComplEx and DistMult, the classical Adamic--Adar
and Common-Neighbours heuristics, a training-free recurrence baseline, and a
ChronoBERT point-in-time-clean semantic baseline, by MRR and Hits@$k$ under a
conservative raw-ranking protocol with a recurring/novel decomposition,
bootstrap confidence intervals and seed replication. Recurring relationships
are highly predictable while genuinely novel ones sit near a low ceiling for
every method; the strongest system is a recurrence$\rightarrow$ChronoBERT
backoff; the method ranking replicates across all three graph constructions;
and the temporal network's advantage over static structure appears only as data
density grows. This qualified null result for temporal structure is reported as
a substantive finding, and a cutoff-controlled ChronoBERT comparison bounds
extractor lookahead bias below measurement resolution ($<0.002$ MRR).

The contributions to science, discussed further in Section~1.4, are as follows.

\textbf{A built, leakage-controlled temporal knowledge graph} in
FinDKG-compatible quadruple format, with themes as CONCEPT and EVENT nodes and a
reproducible point-in-time audit trail, where graph and feature leakage are
prevented by design and extractor hindsight is measured rather than assumed away.

\textbf{A validated walk-forward theme-discovery system}: a definition and
measurement of link-formation velocity and acceleration on a temporal knowledge
graph, a dual-detector burst account of emergence, and a causal
theme-discovery system with data-driven noise controls, validated by lead time
against a frozen gold list and by a placebo control.

\textbf{A decisive, replicated benchmark of edge-level predictability}:
establishing that, under point-in-time discipline, recurrence dominates,
a temporal graph network improves on static structure only as data density
grows, and a recurrence$\rightarrow$ChronoBERT backoff is the strongest system
--- with the method ranking replicated across three graph constructions and
extractor lookahead bias bounded below measurement resolution.

\textbf{An open, reproducible methodology and dataset}: source code and the
derived knowledge-graph dataset (triples and indices only; no article text is
redistributed), released to enable replication and extension.
% =====================================================================
\end{abstract}

\chapter*{Impact Statement}
\addcontentsline{toc}{chapter}{Impact Statement}
The contributions to industry of this dissertation complement its contributions
to science. The method produces a systematic, auditable signal for detecting
emerging investment themes early, before they are widely reported: because every
edge is timestamped by publication date and evaluated under strict point-in-time
discipline, the signal is suitable for production use by quantitative asset
managers, where it can complement existing factor models rather than replace
them. Velocity and burst detection on a temporal financial
knowledge graph extends naturally to risk and compliance applications such as
early detection of contagion in counterparty networks, ESG-controversy
monitoring, and sector-rotation or macro-regime timing. By building the graph
from open data with an open-source extractor and releasing both the code and the
derived knowledge-graph dataset (no source article text is redistributed), the
work lowers the barrier for practitioners and
regulators to reproduce and extend these capabilities, and demonstrates that
open tooling under disciplined point-in-time evaluation can rival closed,
vendor-supplied alternatives. Any direct investment application is treated as a
downstream use of the method rather than its aim.

\tableofcontents
\listoffigures
\listoftables

% ---------------------------------------------------------------------
% MAIN MATTER
% ---------------------------------------------------------------------
\cleardoublepage
\pagenumbering{arabic}
\setcounter{page}{1}

% =====================================================================
\chapter{Introduction}
% Marking area 1 (Background, Aims & Organisation). Introduction written
% in the structure of the Batrinca (2016) UCL thesis.
This chapter introduces the research and its motivation, states the objectives,
describes the three experiments that constitute the dissertation, lists the
contributions to science, and finishes with the structure of the thesis.

\section{Research Motivations}
Investment themes reshape markets, but they are usually recognised only once
they have been widely reported. By the time a theme such as generative AI,
GLP-1 therapeutics or onshoring is named in commentary and research notes, the
attention-based indicators that analysts watch: headline counts, search
interest, and sentiment scores have already reacted, and much of the
informational advantage has dissipated. The relationships among companies,
sectors, countries and themes, however, begin to change before a theme is
named: firms are linked to a new concept, concepts are linked to one another,
and the density of those connections grows. If that structural change can be
observed as it happens, the \emph{rate} at which relationships form may
anticipate what attention-based indicators only confirm after the fact.

Existing financial knowledge graphs are not well suited to answering this
question. They are largely static or transductive, scoring whether a relation
plausibly exists rather than when and how quickly relations form, and they
rarely audit point-in-time correctness, which is the central threat to validity
in any financial study. This dissertation therefore asks a single question:
does the velocity of relationship formation, measured on a graph built under
strict point-in-time discipline, carry information about which relationships
form next?

We measure and characterise the velocity signal by defining it, smoothing it
without looking ahead, and detecting when it bursts; we then use that signal to
discover emerging themes as they form, and separately benchmark how far
individual future links can be predicted at all, against strong baselines and
with honest evaluation.

\section{Research Objectives}
\label{sec:objectives}
The main objective of this research is to determine whether the dynamics of a
financial knowledge graph anticipate the formation of new relationships, under
conditions strict enough that a positive result could be trusted. To this end a
time-indexed knowledge graph is built from public financial news; the velocity
and acceleration of edge formation are defined and measured; emerging themes are
detected as bursts in that signal and validated against known episodes; and the
predictive value of velocity is tested in a temporal link-prediction model
evaluated against strong baselines.

Point-in-time integrity is the central methodological commitment, and the work
distinguishes two kinds of leakage. Graph and feature leakage are prevented by
design: every edge is dated by the article's publication date rather than by the
date of the event it describes, which is the single most important control, together
with chronological train/validation/test splits, structural masking (at time
$t$ the model observes only edges with timestamp $\leq t$), rolling causal
normalisation, a conservative raw-ranking evaluation, within-week
update-after-scoring ordering, and validation-only hyperparameter and threshold
tuning. Extractor hindsight is measured rather than eliminated: the
language-model extractor used to build the graph carries world-knowledge of the
study period, so the residual is quantified through a blinded-versus-unblinded
diagnostic and a cutoff-controlled comparison against chronologically
consistent models, and reported as bounded-and-measured rather than assumed
absent.

We use the ICKG relation extractor; success is measured on graph-native metrics
rather than as an investment signal; and country (GPE) relations, though
present, are sparse owing to the US-centricity of the corpus and are documented
as a limitation.

\section{Research Experiments}
\label{sec:experiments}
The dissertation builds its substrate first, then characterises the velocity
signal, and finally tests its predictive value. It is split into three
experiments:

\begin{enumerate}
\item \textbf{Building the Time-Indexed Knowledge Graph.} The first experiment
constructs the substrate rather than consuming a pre-built graph. The
open-source ICKG extractor (ICKG-v4.2, a LoRA adapter on Qwen-2.5-14B-Instruct
distilled from GPT-4) is run over the FNSPID full-text news corpus to produce
FinDKG-compatible quadruples (subject, relation, object, time) over twelve entity
and fifteen relation types, with themes realised as CONCEPT and EVENT nodes and
every edge timestamped by publication date. The graph is built at two corpus
scales (201{,}600 and 604{,}800 time-stratified articles), and a data-hygiene
audit shows that 10--24\% of corpus articles are verbatim syndication
duplicates whose removal halves the denoised core; the deduplicated core is
adopted as the canonical graph and the smaller corpus retained as a replication
sample. This experiment is deliberately scoped as bounded infrastructure for
the studies that follow.

\item \textbf{Edge-Formation Velocity, Burst Detection, and Walk-Forward Theme
Discovery.} The second experiment contains the dissertation's central
contribution. It formally defines and
measures the rate and acceleration of new-edge formation per node, relation and
theme using trailing-window, past-only Poisson-surprise normalisation; it
detects bursts with three independent detectors (Kleinberg's automaton, the
CUSUM control chart, and Bayesian online changepoint detection) so that
detected episodes are not an artifact of one algorithm; and it clusters co-accelerating entities on their recent
co-occurrence subgraph into themes, tracked week-over-week into lifelines whose
birth week is the date the system would have flagged the theme live. Two
data-driven noise controls (template-article fingerprinting and a habituation
score separating emerging themes from recurring rituals) require no hand
curation. Detections are validated against a frozen ten-event gold list with
event-anchored birth windows, and specificity is established with a
100{,}000-draw placebo control.

\item \textbf{Temporal Link Prediction under Point-in-Time Discipline.} The
third experiment maps what is and is not forecastable at the individual-edge
level, contextualising the theme system. An RE-GCN-style temporal graph network
is evaluated against ComplEx and DistMult static embeddings, Common-Neighbours
and Adamic--Adar heuristics, a training-free recurrence baseline and a
ChronoBERT semantic baseline, using MRR and Hits@$k$ under a conservative
raw-ranking protocol with a recurring/novel decomposition, bootstrap confidence
intervals and seed replication; the same benchmark is run on all three graph
constructions to test replication. A null result is itself a valid finding.
\end{enumerate}

\section{Contributions to Science}
\label{sec:contributions}
The thesis presents the following contributions to science:

\begin{enumerate}
\item \textbf{A built, leakage-controlled temporal knowledge graph, with a
corpus-hygiene audit.} A time-indexed financial knowledge graph constructed
from public news in FinDKG-compatible quadruple format at two corpus scales,
with themes as CONCEPT and EVENT nodes and a reproducible point-in-time audit
trail; graph and feature leakage are prevented by design, extractor hindsight
is measured rather than assumed away, and a syndication-duplication audit
(10--24\% verbatim duplicate articles in FNSPID) quantifies and corrects a
bias affecting any per-ticker news corpus.

\item \textbf{A validated walk-forward theme-discovery system.} A definition
and measurement of link-formation velocity and acceleration on a temporal
knowledge graph; a dual-detector burst account of emergence; and a causal
theme-discovery system (acceleration $\rightarrow$ community clustering
$\rightarrow$ tracked lifelines, with data-driven noise controls) validated by
lead time against a frozen gold list of ten documented episodes and by a
placebo control --- with the finding that lead time reflects the information
structure of the event: anticipatable themes are detected weeks early, shock
events at the news-flow floor.

\item \textbf{A decisive, replicated benchmark of edge-level predictability.}
Establishing that, under point-in-time discipline, recurrence dominates what is
predictable, genuinely novel edges sit near a low ceiling for every method, a
temporal graph network improves on static structure only as data density grows,
and a recurrence$\rightarrow$ChronoBERT backoff is the strongest system ---
with the method ranking replicated across three graph constructions, and
extractor lookahead bias bounded below measurement resolution by a
cutoff-controlled comparison.

\item \textbf{An open, reproducible methodology and dataset.} Source code and
the derived knowledge-graph dataset (no source article text), released to
enable replication and extension.
\end{enumerate}

\section{Thesis Structure}
\label{sec:structure}
Chapter~\ref{ch:background} reviews the relevant literature on financial
knowledge graphs, large-language-model relation extraction, temporal-graph
representation learning, time encoding, burst and change-point detection, and
point-in-time evaluation. Chapter~\ref{ch:exp1} presents Experiment~1, the
construction of the time-indexed knowledge graph, including the extraction
trial. Chapter~\ref{ch:exp2} presents Experiment~2, the definition and
validation of edge-formation velocity, burst detection, and the walk-forward
theme-discovery system. Chapter~\ref{ch:exp3} presents Experiment~3, the
temporal link-prediction benchmark and its replication across graph
constructions. Chapter~\ref{ch:pit} gathers
the point-in-time integrity framework that governs all three experiments.
Chapter~\ref{ch:conclusions} discusses the findings, their limitations and
threats to validity, and directions for future work.

% =====================================================================
\chapter{Background and Literature Review}
\label{ch:background}
% Marking area 1: literature context + extra-curricular reading.
This chapter situates the dissertation within the prior work on which it draws.
The research sits at the intersection of several literatures, and this review
organises them into six strands: the construction of financial knowledge graphs
with large language models, which supplies the substrate; the static
knowledge-graph embeddings that serve as link-prediction baselines; the temporal
knowledge-graph forecasting models, together with the recent recurrence-baseline
critique that tempers expectations of them; the classical link-prediction
heuristics that predate the embedding literature; the detection of bursts and
change-points in event streams, which underpins the emergence analysis; and the
point-in-time integrity and lookahead-bias considerations that govern any
financial study built on a pretrained language model. Each strand is reviewed in
turn and related to the experiments that follow.

\section{Financial Knowledge Graphs and LLM-Based Construction}
A knowledge graph encodes entities and the typed relationships between them. A
temporal, or dynamic, knowledge graph additionally records the time at which each
relationship holds, so that the graph can be reconstructed at any past instant
and its evolution studied over time. This temporal property is central to the
present work, since the dissertation is concerned not with whether a relationship
exists but with how quickly relationships form.

The most directly relevant system in the financial domain is FinDKG, introduced
by \citet{Li2024FinDKG}. FinDKG defines a fixed ontology of twelve entity types
and fifteen relation types, together with a standard temporal knowledge-graph
file format comprising entity and relation index files, a time index, and
chronologically split lists of quadruples. Crucially for this dissertation, the
authors also release the Integrated Contextual Knowledge Graph (ICKG) relation
extractor, an open large-language-model extractor fine-tuned for the task of
knowledge-graph construction. This research adopts the FinDKG ontology and
on-disk format, but it builds its own graph with the ICKG extractor rather than
consuming the released FinDKG graph, so that the choice of corpus, the time
resolution, and the point-in-time discipline all remain under the author's
control.

Building the graph by extraction rather than consuming a pre-built one makes the
dissertation a consumer of several lines of recent language-model research. The
ICKG release employed here is a Low-Rank Adaptation (LoRA) adapter, in the sense
of \citet{Hu2022LoRA}, applied to the instruction-tuned Qwen-2.5 model
\citep{Qwen2024}. LoRA fine-tunes a small set of injected low-rank matrices while
freezing the base weights, which makes a fixed and fully documented extractor
practical to host locally. A local, fixed extractor is preferred to a frontier
model accessed through a commercial application-programming interface (API)
precisely because its knowledge cutoff is controllable and documented, a property
that matters for the point-in-time argument developed in Chapter~\ref{ch:pit}.
High-throughput extraction over the corpus relies on vLLM \citep{Kwon2023vLLM},
whose PagedAttention memory manager enables efficient batched generation on the
available hardware. The corpus itself is FNSPID \citep{Dong2024FNSPID}, which
aligns roughly 15.7 million timestamped financial-news records with stock-price
series for the S\&P~500 universe over the period 1999 to 2023; of these, the
present work uses only the full-text news component.

\section{Knowledge-Graph Embeddings and Link-Prediction Baselines}
Predicting which links a knowledge graph will contain is known as the
link-prediction problem, and the strongest non-temporal baselines for it
originate in the static knowledge-graph-embedding literature. The translational
model TransE, proposed by \citet{Bordes2013TransE}, represents a relation as a
translation between the vector embeddings of its subject and object entities. As
important for the present work as the model itself is the evaluation protocol
that accompanied it: the filtered entity-ranking procedure, reported through Mean
Reciprocal Rank (MRR) and Hits@$k$ over corrupted candidate entities, which has
since become the standard means of assessing link prediction and is adopted
throughout this dissertation --- in its conservative \emph{raw} (unfiltered)
form, so that no evaluation-time knowledge of other true future facts enters a
prediction; this deflates absolute scores relative to the filtered convention
and is the stricter choice under point-in-time discipline. Two later models refine the scoring function.
DistMult \citep{Yang2015DistMult} scores a triple with a bilinear diagonal form,
and ComplEx \citep{Trouillon2016ComplEx} extends this idea to complex-valued
embeddings, whose Hermitian product is able to represent both symmetric and
antisymmetric relations. Because these models score the structural plausibility
of a triple independently of whether it has recurred in the past, they are
informative baselines on genuinely novel edges, and it is in exactly this role
that they are used in Experiment~3.

\section{Temporal Knowledge Graph Forecasting}
Forecasting on a temporal knowledge graph, also called the extrapolation setting,
requires that only history strictly preceding the query time may be used to
predict future links. This setting has been dominated by recurrent neural
architectures. RE-NET, introduced by \citet{Jin2020RENET}, infers graph structure
autoregressively, generating the events of each timestep conditioned on those of
the preceding ones. RE-GCN, due to \citet{Li2021REGCN}, learns time-evolving
representations of entities and relations by passing each snapshot of the graph
through a relational graph convolutional network and folding the result into a
recurrent state, with a decoder that scores candidate future edges. The temporal
graph neural network implemented in this dissertation follows the RE-GCN design.

A recent and directly relevant critique tempers expectations of these models.
\citet{Gastinger2024Recurrence} show that a simple, training-free recurrence
baseline, which scores a candidate fact according to how often that exact fact
has occurred in the past, rivals or outperforms many neural temporal
knowledge-graph forecasters, and they argue that part of the apparent progress in
the field is an artefact of comparison against weak baselines. This finding
directly motivates the careful baselining adopted in Experiment~3, in which the
recurrence baseline is treated not as a foil but as a first-class competitor. It
also anticipates one of the dissertation's own results, namely that the temporal
graph neural network adds little forecasting advantage over static structure on
edges that are genuinely new.

\section{Classical Link-Prediction Heuristics}
Before the embedding and neural methods described above, link prediction was
studied as a purely structural problem in the social-network literature.
\citet{LibenNowell2007} surveyed a family of proximity measures for predicting
which links would later appear in a network, the simplest of which counts the
neighbours that two nodes already share. The index of \citet{AdamicAdar2003}
refines this count by summing over the shared neighbours of a candidate pair,
weighting each by the inverse logarithm of its degree so that rarer common
neighbours count for more. These heuristics encode assumptions of triadic closure
and homophily that hold well in social graphs. Whether the same assumptions
transfer to a financial co-mention graph is an empirical question rather than a
settled one, and a negative answer would itself be informative; for this reason
the Common Neighbours and Adamic-Adar heuristics are retained as baselines in
Experiment~3.

\section{Burst and Event Detection in Streams}
Detecting the moment at which a relationship or a theme suddenly intensifies is a
burst-detection problem. The standard approach is the model of
\citet{Kleinberg2003Bursty}, which represents the stream as the output of an
automaton whose states correspond to different rates of arrival. In its simplest,
two-state form, the optimal state sequence is recovered by trading the improved
likelihood of attributing a stretch of high activity to an elevated-rate state
against a cost incurred each time the automaton switches into that state, so that
only sustained runs of high activity are reported as bursts. The two-state model
tends to merge distinct but adjacent episodes into a single long burst, a
limitation that the hierarchical variant addresses by separating them according
to intensity.

Kleinberg's automaton is not the only lens on the problem. A complementary and
much older tradition treats the same question as \emph{sequential change
detection}: the cumulative-sum (CUSUM) control chart of \citet{Page1954CUSUM}
accumulates standardised deviations from a baseline rate and raises an alarm
when the cumulative excess crosses a decision interval, so that it responds to
sustained shifts in level rather than to discrete spikes. A third treatment is fully Bayesian:
the Bayesian online changepoint detection of \citet{AdamsMacKay2007BOCD}
maintains an exact online posterior over the time since the last changepoint
(here with a conjugate Gamma--Poisson model; changepoints are back-dated to
the inferred onset).
This dissertation runs all three detectors over the edge-formation series, so
that detected episodes are demonstrably not an artifact of a single algorithm:
empirically they agree on where activity is elevated while segmenting it at
different granularities (Kleinberg isolating discrete events, CUSUM and BOCD
sustained regimes). Since interval widths differ, recall is judged against
each detector's own chance level (its dilated flagged-week coverage --- the
recall it would achieve on random event weeks): on the gold events recall runs
78--94\% against chance of 16--29\% (an excess of +59 to +66 points), with a
three-detector union of 18/18.

Detecting that individual entities accelerate is, however, not yet theme
discovery: an emerging theme is a \emph{group} of related entities accelerating
together. For this step the dissertation clusters co-accelerating entities on
their recent co-occurrence subgraph using the Louvain community-detection
algorithm of \citet{Blondel2008Louvain}, a fast greedy optimiser of modularity,
and tracks the resulting communities week over week into theme lifelines ---
a construction in the spirit of the dynamic-community-tracking literature.
The resulting event catalogue and theme lifelines are validated against an
independently compiled list of known episodes.

\section{Point-in-Time Integrity and Lookahead Bias}
The central methodological risk in any financial study that relies on a model
trained on data from after the prediction date is lookahead bias, that is, the
leakage of future information into a prediction that is supposed to depend only on
the past. The risk is acute here because the extractor used to build the graph is
a pretrained language model whose knowledge cutoff postdates much of the corpus,
so the resulting hindsight must be measured rather than simply assumed away. The
most applicable work is that of \citet{He2025Chrono}, who train chronologically
consistent language models, including ChronoBERT, exclusively on text available
up to a stated date, so that such models are free of lookahead bias by
construction. Comparing two versions of the same model that differ only in their
training cutoff yields a capacity-controlled estimate of the bias, which the
authors report to be modest. These models descend architecturally from BERT
\citep{Devlin2019BERT}, the bidirectional transformer encoder. Two alternatives
sometimes suggested for time-aware financial language modelling are worth
distinguishing here. FinBERT \citep{Araci2019FinBERT}, a sentiment-oriented
BERT fine-tune, applies no chronological control to its training text and is
therefore itself exposed to the lookahead risk under study, which is why the
chronologically consistent family is preferred in this dissertation. And the
sinusoidal positional encodings of the Transformer \citep{Vaswani2017Attention}
encode token position within a sequence rather than calendar-time information
availability; point-in-time correctness is accordingly enforced here at the
data and evaluation-protocol level, where it is verifiable, rather than inside
the encoder. Experiment~3 uses
ChronoBERT entity embeddings in two ways: as a point-in-time-clean semantic
baseline in its own right, and, through the cutoff-controlled comparison, as the
instrument for the lookahead-bias diagnostic developed in Chapter~\ref{ch:pit}.
Throughout the dissertation, uncertainty in the reported metrics is quantified
using the bootstrap \citep{Efron1979Bootstrap}, which resamples the test
instances with replacement in order to estimate the sampling distribution of a
statistic without recourse to parametric assumptions.

% =====================================================================
\chapter{Experiment 1: Building the Time-Indexed Knowledge Graph}
\label{ch:exp1}
% Marking area 2: the build is a substantial, bounded software deliverable.
\section{Introduction}
This experiment builds the substrate on which the rest of the dissertation
stands: a time-indexed financial knowledge graph in which every edge can be
traced to a dated public news article. It is deliberately scoped as bounded
infrastructure. Correctness, point-in-time integrity and reproducibility were
prioritised over modelling sophistication, because any flaw introduced here is
inherited by every result in Chapters~\ref{ch:exp2} and \ref{ch:exp3}. The
experiment delivers three versioned graph constructions, a data-hygiene audit
of the source corpus, and the fixed chronological splits used throughout.

\section{Data: the FNSPID Corpus}
FNSPID \citep{Dong2024FNSPID} aligns roughly 15.7 million timestamped
financial-news records with stock-price series for the S\&P~500 universe; only
the full-text news component is used here. The study window is 2017--2023 (84
months, 365 weekly buckets), chosen to cover the COVID-19 shock and the
generative-AI surge together with a multi-year pre-episode baseline.

Raw article counts in FNSPID are strongly skewed towards recent years, so a
graph built from a uniform random sample would confound structural change with
sampling density. Subsets are therefore \emph{time-stratified}: a single
streaming pass over the corpus performs per-month reservoir sampling
\citep{Vitter1985Reservoir} with a fixed seed, keeping only articles whose
bodies exceed 200 words. Two scales are drawn: 2{,}400 articles per month
($=201{,}600$) and 7{,}200 per month ($=604{,}800$). Even monthly coverage
guarantees that week-on-week changes in the graph reflect changes in reported
structure rather than changes in corpus volume.

\section{Extraction with ICKG-v4.2}
Each article is passed to ICKG-v4.2, the FinDKG relation extractor
\citep{Li2024FinDKG}, hosted locally via vLLM \citep{Kwon2023vLLM} with greedy
decoding and the verbatim FinDKG prompt. The extractor returns typed triples
over the fixed ontology of twelve entity types and fifteen relation types,
with themes realised as CONCEPT and EVENT nodes.

A practical finding worth recording: the model emits triples in two formats
interchangeably --- JSON arrays of arrays ($\approx$65\% of articles) and
Python-style parenthesised tuples ($\approx$35\%) --- and a parser tolerant of
both (plus a filter for purely numeric ``entities'') raised the usable-output
rate in the pilot from 35\% to over 99\%. The full runs achieved 0\% malformed
generations at both scales, $\sim$18.5 valid triples per article, and 11.5
million raw triples at the larger scale. A precision spot-check found both
entities of a triple verbatim in the source text for 83\% of a sampled set ---
a lower bound, since the extractor frequently simplifies surface forms.

\section{Edge Construction and Publication-Date Timestamping}
Every edge is timestamped by the article's \emph{publication date}, never by
the date of any event the text describes; this is the single most important
leakage control in the pipeline. Triples are aggregated into FinDKG-format
quadruples with weekly time-bucketing, where the weight of an edge-week is the
number of supporting mentions in that week. Chronological splits are fixed
once, by week: train 0--254, validation 255--309, test 310--364 (a 70/15/15
division of the 365 weeks).

\section{Compute and Throughput}
Extraction ran on a UCL cluster node with two RTX Pro 6000 GPUs, one vLLM
instance per GPU over interleaved article shards, taking $\sim$14\,h for the
200k corpus and $\sim$43\,h for the 600k corpus, with CSV-append checkpointing
so that interrupted runs resume without repeated work. Every subsequent
analysis in this dissertation --- including all of Chapters~\ref{ch:exp2} and
\ref{ch:exp3} --- runs on a single consumer GPU with 6\,GB of memory, which is
a deliberate reproducibility property of the design.

\section{The Syndication-Duplication Audit and the Deduplicated Core}
FNSPID stores one row per article--ticker pair, so an article syndicated to
$k$ tickers can enter a subset up to $k$ times verbatim. Hashing article
bodies reveals that \textbf{10.5\%} of the 200k subset and \textbf{24.3\%} of
the 600k subset are exact duplicates; of 102{,}426 duplicate groups in the
larger subset, 101{,}346 differ only in the ticker column, and only 476 span
more than one week. The damage is therefore concentrated in \emph{same-week
weight inflation}: a threshold such as ``keep edges supported by at least
three mentions'' can be satisfied by a single syndicated article, so
duplication silently biases which edges enter a denoised core.

Removing duplicates (keeping one copy per body) and rebuilding with unchanged
thresholds halves the denoised 600k core, from 41{,}888 to 21{,}216 entities
(212{,}294 edge-weeks). This \emph{deduplicated core} is the only
construction whose thresholds mean what they claim --- support by $N$
independent articles --- and it is adopted as the canonical graph for the
remainder of the dissertation. All principal conclusions are re-verified on
it, and the audit itself is reported as a finding that applies to any
per-ticker news corpus.

\section{Chronological Splits and Versioning}
The on-disk format permits exact reconstruction of \Gle{} for any week $t$.
Three constructions are versioned separately and summarised in
Table~\ref{tab:constructions}: the 200k build with a weight~$\geq 2$ /
degree~$\geq 2$ core, the 600k build with a stricter weight~$\geq 3$ /
degree~$\geq 3$ core, and the deduplicated 600k core. The 200k build is
retained as a replication sample: Chapter~\ref{ch:exp3} runs its full
benchmark on all three.

\begin{table}[ht]
\centering
\caption{Graph constructions. Novelty is the share of test edges whose exact
triple never appeared before its prediction week.}
\label{tab:constructions}
\begin{tabular}{lrrrr}
\toprule
Construction & Articles & Entities & Edge-weeks & Test novelty \\
\midrule
200k full graph & 201{,}600 & 628{,}765 & 2{,}379{,}438 & --- \\
200k core ($w\geq2$, $d\geq2$) & & 46{,}026 & 307{,}647 & 72.4\% \\
600k full graph & 604{,}800 & 1{,}298{,}833 & 6{,}126{,}175 & --- \\
600k core ($w\geq3$, $d\geq3$) & & 41{,}888 & 493{,}238 & 64.0\% \\
\textbf{600k deduplicated core (canonical)} & 457{,}619 & \textbf{21{,}216}
& \textbf{212{,}294} & \textbf{61.6\%} \\
\bottomrule
\end{tabular}
\end{table}

\section{Extraction-Noise Controls and Limitations}
Noise is measured rather than hidden. In the full graphs roughly two thirds of
entities are singletons, motivating the degree/weight cores. Templated
financial-media boilerplate is measured with the slot-fingerprint method of
Section~\ref{sec:noisecontrols} and removed by an opt-in filter (on the
canonical core: 291 entities carrying 15.6\% of edge-instances). Only light
surface-form normalisation is applied; entity resolution is deliberately out
of scope, but a crude corporate-suffix merge bounds its effect --- test
novelty falls from 72.4\% to 68.4\% (200k) and from 64.0\% to 60.1\% (600k),
so fragmentation is real but secondary. Country (GPE) relations are sparse
owing to the US-centric corpus and are documented as a limitation rather than
hand-curated.

\section{Results and Discussion}
The substrate meets its three design goals. It is \emph{reproducible}: fixed
seeds, checkpointed extraction, and a fully scripted path from raw corpus to
splits. It is \emph{point-in-time clean by construction} at the graph level:
publication-date timestamps and chronological splits are enforced before any
model sees the data. And it is \emph{honest about its noise}: duplication,
fragmentation and boilerplate are quantified, with the deduplicated core as
the corrected canonical object. Scale behaves as hoped --- tripling the
corpus deepens evidence per relationship (test novelty falls monotonically
across the constructions in Table~\ref{tab:constructions}) --- and the
duplication audit stands as a transferable methodological result about
building graphs from per-ticker news data.

% =====================================================================
\chapter{Experiment 2: Edge-Formation Velocity, Burst Detection, and Walk-Forward Theme Discovery}
\label{ch:exp2}
% Core contribution + validation.
\section{Introduction}
This chapter contains the dissertation's central contribution: a walk-forward
system that discovers emerging themes as they form. The pipeline has four
stages --- score each entity's formation acceleration causally, detect bursts,
cluster co-accelerating entities on the graph into themes, and track themes
into lifelines whose birth week is the date the system would have flagged them
live --- followed by validation against a frozen gold list with a placebo
control. All analyses run on the canonical deduplicated core of
Chapter~\ref{ch:exp1}.

\section{Velocity and Acceleration: Definitions}
A \emph{formation} is the first week in which a distinct
(subject, relation, object) triple appears. For any entity, relation or theme,
let $c_t$ denote its formation count in week $t$. Velocity and acceleration
are the first and second differences,
\begin{equation}
v_t = \Delta c_t = c_t - c_{t-1}, \qquad
a_t = \Delta^2 c_t = v_t - v_{t-1},
\label{eq:veldef}
\end{equation}
per node, per relation and per theme. Two window conventions serve two
purposes: \emph{centred} windows for descriptive reporting, where a peak
should align with the event that caused it, and \emph{strictly trailing}
windows for the detection system, which must be causal to claim real-time
operation.

The raw second difference in Equation~\ref{eq:veldef} is an unbiased but
extremely noisy acceleration estimate for sparse count series. The detector
therefore uses its windowed, variance-stabilised form --- the surprise
statistic of the next section, which contrasts the recent-window formation
rate with the entity's own baseline rate. This is the operational form of
``anomalous acceleration'' throughout the dissertation: a smoothed estimate of
the velocity change, robust to the count noise that dominates $\Delta^2 c_t$.

\section{Causal Acceleration Scoring}
At week $t$, with a recent window of $r=4$ weeks and a baseline window of
$b=26$ weeks, define for every entity
\begin{equation}
O_t = \sum_{s=t-r+1}^{t} c_s, \qquad
E_t = \frac{r}{b} \sum_{s=t-r-b+1}^{t-r} c_s, \qquad
z_t = \frac{O_t - E_t}{\sqrt{E_t + 1}} .
\label{eq:surprise}
\end{equation}
Both windows are strictly trailing. Under a Poisson null the denominator
stabilises the variance; the $+1$ floor keeps $z_t$ finite for zero baselines,
so brand-new entities (``Coronavirus'' in January 2020) score high by design
--- new-entity emergence is a feature, not an artifact.

Weekly formation counts are in fact strongly over-dispersed. A
negative-binomial variant replaces the denominator with
$\sqrt{E_t + \alpha E_t^2 + 1}$, where the dispersion $\alpha = 3.75$ is
estimated once by the method of moments on training weeks only. Raising the
variance raises the detection bar: the Poisson floor is the high-recall
configuration and the NB correction the high-precision one, and the gold-list
conclusion of Section~\ref{sec:themevalidation} holds under both. Two
disclosures apply: $\alpha$ is estimated on weeks 0--254 but applied at every
$t$, which is conservative for earlier events; and the template blacklist of
Section~\ref{sec:noisecontrols} is a static full-period list --- the one
non-trailing input to an otherwise walk-forward system.

\section{Burst and Change-Point Methods}
Three detectors with different inductive biases are run on the formation
series, so that detected episodes are demonstrably not an artifact of one
algorithm. Kleinberg's two-state automaton \citep{Kleinberg2003Bursty} finds
discrete bursts by a Viterbi trade-off between emission likelihood and a
transition cost for entering the elevated state. The one-sided CUSUM chart
\citep{Page1954CUSUM} accumulates standardised deviations,
$C_t = \max(0,\, C_{t-1} + z_t - k)$ with reference value $k=0.5$, and alarms
when $C_t$ exceeds a decision interval $h=5$, responding to sustained shifts.
Bayesian online changepoint detection \citep{AdamsMacKay2007BOCD} maintains
the exact run-length posterior under a Gamma--Poisson model with hazard
$1/100$; changepoints are back-dated to the inferred onset $t - \ell_t$, where
$\ell_t$ is the MAP run length.

Because the three flag intervals of very different widths --- Kleinberg
isolates events, CUSUM and BOCD sustained regimes (their global intervals
agree at Jaccard 0.81) --- recall must be judged against each detector's own
chance level, defined as its dilated flagged-week coverage: the recall it
would achieve on random event weeks. On an 18-event gold list of documented
2017--2023 episodes (Table~\ref{tab:detectors}), all three detect far above
chance with comparable excesses, and their union recovers every event.
These detectors are retrospective descriptive tools (their baselines use the
full series); the walk-forward system is the theme detector itself.

\begin{table}[ht]
\centering
\caption{Entity-level event recovery on the canonical core: 18 documented
episodes, hit $=$ a burst interval overlapping the event week $\pm 3$ weeks.}
\label{tab:detectors}
\begin{tabular}{lrrr}
\toprule
Detector & Recall & Chance (dilated coverage) & Excess \\
\midrule
Kleinberg & 14/18 $=$ 78\% & 16\% & $+62$pp \\
CUSUM & 15/18 $=$ 83\% & 24\% & $+59$pp \\
BOCD & 17/18 $=$ 94\% & 29\% & $+66$pp \\
Union & \textbf{18/18 $=$ 100\%} & --- & --- \\
\bottomrule
\end{tabular}
\end{table}

\section{Theme Clustering and Lifelines}
Individual acceleration is not yet a theme: an emerging theme is a group of
related entities accelerating together. Each week, entities passing the gate
($z_t \geq 3$ and $O_t \geq 6$) are placed on their trailing eight-week
co-occurrence subgraph and clustered with Louvain
\citep{Blondel2008Louvain} (fixed seed; clusters of fewer than three members
discarded). Clusters are ranked by total excess formations
$\sum_i (O_t - E_t)_i$ and the top eight per week retained. Themes are
tracked week-over-week into \emph{lifelines} by member overlap (Jaccard
$\geq 0.3$), matching live lifelines first --- a continuation beats a
re-activation regardless of overlap --- and allowing a lifeline to be claimed
by at most one cluster per week; a re-activation after a gap of more than
eight weeks rejoins its old lifeline rather than counting as a new emergence.
A lifeline's \emph{birth week} is the date the system would have flagged the
theme live.

\section{Data-Driven Noise Controls}
\label{sec:noisecontrols}
Two failure modes require control, and both controls are measured from data
rather than hand-curated. First, \emph{templated screener articles} (options
screens, quantitative stock reports) inject formulaic entities whose ramp-ups
mimic themes. Slot-normalised fingerprints of titles and body openings ---
rare tokens replaced by a placeholder so that ``company-name slots'' collapse
--- identify 26.4\% of articles as template-generated; the 2{,}010 entities
drawing $\geq 60\%$ of their mentions from such articles are excluded from
candidacy. Second, \emph{recurring rituals} (earnings-season screens, index
complexes) accelerate every quarter without being new. A \emph{habituation}
score classifies a newborn lifeline as EMERGING only if the excess-weighted
share of members with fewer than two prior acceleration episodes is at least
one half --- so a theme is emerging only when the entities driving it are new
phenomena. A stricter configuration additionally excludes walk-forward
ubiquitous entities (active in more than half of all past weeks); it
dissolves star-shaped screener clusters at some cost in lead time, and both
configurations are reported.

\section{Validation: Frozen Gold List, Lead Times, and Placebo Control}
\label{sec:themevalidation}
Validation uses a frozen list of ten externally dated 2020--2023 episodes,
word-boundary entity matching, and an event-anchored birth window ($-13$ to
$+8$ weeks) so that pre-existing topical lifelines cannot masquerade as
detections. Table~\ref{tab:goldlist} gives the per-event results for the
default configuration.

\begin{table}[ht]
\centering
\caption{Gold-list lead times (default configuration). Lead $=$ birth week
$-$ event week; negative is early.}
\label{tab:goldlist}
\begin{tabular}{llll}
\toprule
Gold theme & Event week & Birth week & Lead \\
\midrule
COVID market crisis & 2020-02-24 & 2020-02-03 & $-3$ wk \\
Oil crash (negative WTI) & 2020-04-20 & 2020-04-06 & $-2$ wk \\
Vaccine race (Warp Speed) & 2020-05-15 & 2020-03-16 & $-8$ wk \\
Vaccine efficacy readout & 2020-11-09 & 2020-09-28 & $-6$ wk \\
GameStop squeeze & 2021-01-25 & 2021-01-25 & $+0$ wk \\
Russia--Ukraine war & 2022-02-21 & 2022-02-28 & $+1$ wk \\
FTX collapse & 2022-11-07 & --- & miss \\
AI / ChatGPT wave & 2022-11-28 & --- & miss \\
SVB banking crisis & 2023-03-06 & 2023-03-13 & $+1$ wk \\
UAW strike & 2023-09-11 & 2023-09-18 & $+1$ wk \\
\midrule
\multicolumn{4}{l}{\textbf{8/10 detected; median lead $-1$ week; 62\%
early-or-on-time.}} \\
\bottomrule
\end{tabular}
\end{table}

Specificity is established by a placebo control: applying the identical
detection criterion to random event weeks (100{,}000 Monte-Carlo draws over
the per-theme chance rates) yields an expected $\approx 2.2/10$; the observed
8/10 has $p \approx 10^{-4}$ (strict configuration: expected $1.9/10$,
$p = 2\times10^{-5}$). Robustness: every cell of the $3\times3$ window grid
(recent $\in \{2,4,8\}$, baseline $\in \{13,26,52\}$ weeks) detects 7--8/10
with placebo $p < 4\times10^{-4}$, shorter windows detecting slightly later
and longer slightly earlier (full grid in Appendix~B); and the
negative-binomial configuration detects 7/10 at a chance rate of $0.9/10$.

\section{Results and Discussion}
The central empirical pattern is that \emph{lead time reflects the
information structure of the event itself}. Anticipatable themes --- ones
whose build-up is reported before the market-moving moment --- are detected
weeks early: the COVID market crisis three weeks before the crash, the
vaccine race eight weeks before Operation Warp Speed was announced, the
efficacy theme six weeks before the Pfizer readout. Pure shocks --- an
invasion, a bank run, a strike --- land at the news-flow floor of $+1$ week,
because no public reporting precedes them. The two misses are coverage-honest
rather than algorithmic: FTX and OpenAI are barely present in a Nasdaq-equity
news corpus. One label is worth noting: GameStop is detected in the squeeze
week but classified \emph{recurring}, since squeeze-adjacent entities had
prior acceleration episodes --- an honest reflection of the habituation
bookkeeping rather than an error. Beyond the gold list the system surfaces
correctly dated discoveries including the 2019 cannabis boom
(HEXO/Aurora/Tilray, born 2019-03-18), the EV/SPAC wave (Nikola/Nio, born
2020-06-22) and the stay-at-home complex (Zoom/Kroger/Virgin Galactic, born
2020-04-06).

% =====================================================================
\chapter{Experiment 3: Temporal Link Prediction under Point-in-Time Discipline}
\label{ch:exp3}
% Marking areas 2 & 4: ML benchmark + critical analysis; contextualises Exp 2.
\section{Introduction}
Where Chapter~\ref{ch:exp2} asks \emph{what is emerging}, this chapter asks
how far the graph's future can be predicted at all at the level of individual
edges, contextualising the theme system. The benchmark deliberately follows
the recurrence-baseline critique of \citet{Gastinger2024Recurrence}: the
training-free baseline is treated as a first-class competitor, uncertainty is
quantified, and the same benchmark is repeated on all three graph
constructions as a replication test.

\section{Task and Protocol}
The task is temporal knowledge-graph forecasting (extrapolation): for each
test quadruple $(s, r, o, t)$, rank all candidate entities for $(s, r, ?, t)$
and $(?, r, o, t)$ using only history strictly before week $t$. Performance
is Mean Reciprocal Rank and Hits@$k$,
\begin{equation}
\mathrm{MRR} = \frac{1}{|Q|} \sum_{q \in Q} \frac{1}{\mathrm{rank}_q},
\qquad
\mathrm{Hits@}k = \frac{1}{|Q|} \sum_{q \in Q}
\mathbb{1}[\mathrm{rank}_q \leq k],
\end{equation}
averaged over both directions, under the conservative \emph{raw} ranking of
Section~2.2 and with mean-rank tie handling (ties share the expected rank, so
count-based baselines are neither inflated nor penalised). Every result is
decomposed by a key property of the data: a test edge is \emph{recurring} if
its exact triple appeared before its prediction week and \emph{novel}
otherwise. On the canonical core 61.6\% of test directions are novel.

\section{Models}
The temporal model is an RE-GCN-style network \citep{Li2021REGCN}: a
relational graph convolution over each weekly snapshot folded into a
recurrent state, with entity representations re-evolved from their base
embeddings over a trailing window at each prediction time and a decoder
scoring candidate edges. Hyperparameters are selected by validation-based
early stopping on each core's own validation weeks (epoch 22 on the canonical
core), never transplanted across constructions.

\section{Baselines}
Five baseline families bracket the model. The training-free
\emph{recurrence} baseline scores a candidate by how often the exact triple
occurred in prior history \citep{Gastinger2024Recurrence}; \emph{popularity}
ranks by global historical frequency. Static embeddings --- ComplEx
\citep{Trouillon2016ComplEx} and DistMult \citep{Yang2015DistMult} --- are
trained on all pre-test weeks and score structural plausibility independent
of recurrence. Classical heuristics --- Common Neighbours and Adamic--Adar
\citep{AdamicAdar2003, LibenNowell2007} --- operate on a trailing twelve-week
subgraph, with the query entity excluded from its own candidate set. Finally,
ChronoBERT entity-name embeddings \citep{He2025Chrono} with a lightweight
trained scorer provide a point-in-time-clean \emph{semantic} baseline.
Backoff hybrids use the recurrence rank for recurring candidates and a second
scorer for novel ones.

\section{Evaluation Protocol and Uncertainty}
All scorers observe strictly pre-$t$ history, with within-week
update-after-scoring ordering so no same-week information leaks. Confidence
intervals are 95\% paired bootstrap (1{,}000 resamples over test directions),
and --- a methodological lesson of this dissertation --- small gaps are
additionally checked against training-seed replication, because query
resampling alone conditions on a single training run and understates
uncertainty (Chapter~\ref{ch:pit}). The full benchmark is run on the 200k
core, the 600k core and the canonical deduplicated core.

\section{Central Question: Does Temporal Structure Help?}
Table~\ref{tab:benchmark} reports the canonical-core benchmark.

\begin{table}[ht]
\centering
\caption{Link-prediction benchmark on the canonical deduplicated core
(raw ranking; 95\% bootstrap CIs in brackets where computed).}
\label{tab:benchmark}
\begin{tabular}{lrrrr}
\toprule
Method & MRR all & MRR recurring & MRR novel & Hits@10 \\
\midrule
Recurrence & 0.196 & 0.511 & 0.000 & 0.296 \\
Popularity & 0.053 & 0.097 & 0.025 & 0.119 \\
Common Neighbours / Adamic--Adar & 0.025 / 0.029 & 0.054 / 0.062 &
0.008 & 0.060 \\
ComplEx & 0.099 & 0.218 & 0.025 [.024, .026] & 0.193 \\
RE-GCN (early-stopped) & 0.117 & 0.257 & 0.031 & 0.196 \\
ChronoBERT-2022 & 0.105 & 0.206 & 0.035 [.034, .037] & 0.186 \\
Backoff (rec $\rightarrow$ RE-GCN) & 0.215 & 0.511 & 0.031 & 0.337 \\
\textbf{Backoff (rec $\rightarrow$ ChronoBERT)} & \textbf{0.218 [.215,
.221]} & 0.511 & 0.035 & \textbf{0.342} \\
\bottomrule
\end{tabular}
\end{table}

Four findings follow. \emph{Recurrence dominates what is predictable}:
recurring edges reach MRR $\approx 0.5$ while every method sits near a low
ceiling on novel edges. \emph{Semantic beats structural on novel edges}: the
ChronoBERT baseline exceeds ComplEx by $+0.011$ MRR
[95\% CI $+0.009, +0.013$], a gap an order of magnitude above both seed noise
and the lookahead bound, making the recurrence$\rightarrow$ChronoBERT backoff
the strongest system. \emph{Temporal structure helps only with density}: on
the sparse 200k core the temporal network merely matches ComplEx (novel MRR
0.025 for both), and its score is invariant to shuffling the order of history
--- but on the denser cores a modest advantage emerges (0.031 vs 0.025 novel
on the canonical core). This qualified null is reported as a substantive
finding. \emph{The ranking replicates}: the method ordering ---
recurrence-backoff best overall; ChronoBERT $>$ RE-GCN $\geq$ ComplEx $\gg$
heuristics on novel edges --- is identical on all three constructions
(Appendix~B), across corpus scale, thresholds and deduplication.

\section{Results and Discussion}
The picture is bimodal: the past is legible, the future is genuinely
surprising. What recurs can be predicted almost for free, novel formations
resist every model family tested, and the best available lever on novel edges
is semantic knowledge of \emph{what an entity is} rather than temporal or
structural pattern. Cross-construction comparisons of absolute MRR are
avoided throughout --- raw ranks over candidate sets of 21k--46k entities are
not commensurable --- so the defensible cross-corpus claims are the novelty
trajectory (72.4\% $\rightarrow$ 64.0\% $\rightarrow$ 61.6\%) and the
invariant method ranking. The failure of triadic-closure heuristics
(novel MRR $\approx 0.008$) is itself informative: the assumptions that hold
in social networks transfer poorly to a financial co-mention graph. These
results frame the theme system of Chapter~\ref{ch:exp2}: individual novel
edges are near-unforecastable, yet their aggregate acceleration is exactly
where exploitable signal lives.

% =====================================================================
\chapter{Point-in-Time Integrity}
\label{ch:pit}
% Marking area 4: the methodological backbone governing all three experiments.
\section{Two Kinds of Leakage}
Two distinct threats are conflated by the umbrella term ``lookahead bias'',
and this dissertation treats them separately. \emph{Graph and feature
leakage} is the use, at prediction time $t$, of any data dated after $t$; it
is a property of the pipeline and can be eliminated by construction.
\emph{Extractor hindsight} is subtler: the language model that built the
graph was trained on text postdating the corpus, so its extractions could in
principle encode knowledge of outcomes. It cannot be eliminated without
retraining the extractor; it can only be measured and bounded.

\section{Prevented by Design}
The first threat is addressed by mechanism, not by care: publication-date
timestamping of every edge (the single most important control); chronological
train/validation/test splits fixed once; structural masking, so any model at
time $t$ observes only \Gle; strictly trailing windows and rolling causal
normalisation in the detection system; conservative raw ranking, so no
evaluation-time knowledge of other true future facts enters a prediction;
within-week update-after-scoring ordering; and validation-only tuning. Two
disclosed exceptions to the trailing-statistics rule are the static template
blacklist and the train-estimated NB dispersion of Chapter~\ref{ch:exp2},
both conservative in direction.

\section{Measured, Not Eliminated}
ICKG is distilled from GPT-4 and is not point-in-time clean. Residual
hindsight is quantified with four independent checks. (i)
\emph{Apples-to-apples retraining}: the temporal network trained on identical
data matches static ComplEx, so no comparison rests on an uneven split. (ii)
\emph{Time-shuffle placebo}: the temporal model's score is invariant to
randomising the order of its history snapshots (MRR $0.090 \rightarrow
0.091$), so the null result is intrinsic to the data rather than a tuning
artifact. (iii) \emph{Blinded re-extraction}: re-extracting a 300-article
sample with entity names and dates masked retains the relational structure
(relation-mix total-variation distance 0.10), confirming text-grounded
extraction; only the causal ``impact'' relations lose share ($27.5\%
\rightarrow 21.1\%$), which flags impact-relation analyses --- not link
prediction --- as entity-dependent. (iv) \emph{Cutoff-controlled ChronoBERT
comparison} \citep{He2025Chrono}: scoring the same task with two versions of
the same model differing only in training cutoff (2022, blind to the test
year, vs 2024) bounds lookahead below measurement resolution. The apparent
gap ($+0.0010$ MRR at 200k) initially tested significant under the query
bootstrap, but training-seed replication showed run-to-run noise of the same
size, and on the canonical core the gap is not significant even before that
correction ($+0.0005$, 95\% CI $[-0.0009, +0.0019]$). The honest conclusion
is a bound, not a detection: extractor lookahead is below $\sim$0.002 MRR,
an order of magnitude smaller than the semantic effect it could be confused
with. The methodological lesson --- query resampling conditions on a single
training run and can certify noise as signal --- is reported as a
contribution in its own right.

% =====================================================================
\chapter{Discussion and Conclusions}
\label{ch:conclusions}
% Marking area 4: critical analysis, weaknesses, extensions.
\section{Summary of Findings}
Five findings summarise the dissertation. First, a reproducible,
point-in-time-clean temporal knowledge graph can be built from public news
with an open, locally hosted extractor at practical cost, and its dominant
data-hygiene hazard --- syndication duplication in per-ticker corpora --- can
be measured and corrected. Second, link-formation velocity, causally scored
and clustered on the graph, supports a walk-forward theme-discovery system
that detects 8/10 major 2020--2023 themes at a median lead of $-1$ week,
far above a placebo chance of $\approx 2/10$ ($p \approx 10^{-4}$), robustly
across detector families, dispersion assumptions and window choices. Third,
lead time reflects the information structure of the event: anticipatable
themes are flagged weeks early, pure shocks at the news-flow floor. Fourth,
edge-level forecasting is bimodal --- recurring relationships are highly
predictable, novel ones sit near a low ceiling for every method, temporal
structure helps only as data density grows, and semantic priors are the best
lever on novel edges. Fifth, extractor lookahead bias is bounded below
measurement resolution by a cutoff-controlled comparison.

\section{Strengths and Weaknesses}
The work's principal strength is the validation discipline around its central
claim: a frozen gold list with event-anchored windows, a placebo null, three
burst detectors scored against their own chance levels, dispersion and window
sensitivity, and replication of the benchmark across three graph
constructions. Its corresponding weakness is breadth over depth in places:
the theme system's precision is established at the week level (placebo) but a
lifeline-level false-discovery audit remains manual future work; and the
velocity signal is used for discovery, not yet injected as a feature into the
temporal network, so the feature-ablation form of the velocity question
remains open.

\section{Limitations and Threats to Validity}
\emph{Residual extractor hindsight} is bounded, not zero, and causal impact
relations carry measurable entity-dependence. \emph{Corpus coverage}: a
Nasdaq-equity news corpus under-reports crypto and private-company themes,
which is why FTX and ChatGPT are missed --- detection is bounded by what the
corpus reports. \emph{Duplication}: exact-hash deduplication is a lower
bound; near-duplicates survive. \emph{GPE sparsity}: country relations are
thin in a US-centric corpus. \emph{Metric comparability}: absolute raw MRRs
are not comparable across candidate-set sizes, so cross-construction claims
are restricted to rankings and novelty rates. \emph{Regime dependence}: the
gold list is concentrated in 2020--2023; earlier, calmer regimes contribute
fewer validation events. \emph{Multiple testing}: the gold list and windows
were frozen before evaluation, but ten events remain a small sample.
\emph{Static inputs}: the template blacklist and the NB dispersion are
full-period or train-period statistics, disclosed in
Chapter~\ref{ch:exp2}.

\section{Future Work}
Four extensions follow directly. Injecting velocity and acceleration as
explicit features into the temporal network would answer the feature-ablation
form of the velocity question. Entity resolution --- shown by the crude-merge
bound to move novelty by roughly four percentage points --- is the next
data-quality lever. The detection layer would benefit from a rolling template
blacklist, hierarchical burst detection, and a lifeline-level false-discovery
audit. Finally, the FNSPID price series enable a downstream study of whether
theme leads translate into investable signal; that application is treated
throughout as a use of the method, not its aim.

% =====================================================================
% BIBLIOGRAPHY  (Harvard / author-year; verify all details before submission)
% =====================================================================
\begin{thebibliography}{Liben-Nowell and Kleinberg, 2007}

\bibitem[Adamic \& Adar(2003)]{AdamicAdar2003}
Adamic, L. A. \& Adar, E., 2003. Friends and Neighbors on the Web. {\em Social
Networks}, 25(3), pp. 211--230.

\bibitem[Adams \& MacKay(2007)]{AdamsMacKay2007BOCD}
Adams, R. P. \& MacKay, D. J. C., 2007. Bayesian Online Changepoint Detection.
{\em arXiv preprint} arXiv:0710.3742.

\bibitem[Araci(2019)]{Araci2019FinBERT}
Araci, D., 2019. FinBERT: Financial Sentiment Analysis with Pre-trained
Language Models. {\em arXiv preprint} arXiv:1908.10063.

\bibitem[Blondel et al.(2008)]{Blondel2008Louvain}
Blondel, V. D., Guillaume, J.-L., Lambiotte, R. \& Lefebvre, E., 2008. Fast
Unfolding of Communities in Large Networks. {\em Journal of Statistical
Mechanics: Theory and Experiment}, P10008. arXiv:0803.0476.

\bibitem[Bordes et al.(2013)]{Bordes2013TransE}
Bordes, A. et al., 2013. Translating Embeddings for Modeling Multi-relational
Data. {\em Advances in Neural Information Processing Systems 26 (NeurIPS 2013)}.

\bibitem[Devlin et al.(2019)]{Devlin2019BERT}
Devlin, J., Chang, M.-W., Lee, K. \& Toutanova, K., 2019. BERT: Pre-training of
Deep Bidirectional Transformers for Language Understanding. {\em Proceedings of
NAACL-HLT}. arXiv:1810.04805.

\bibitem[Dong et al.(2024)]{Dong2024FNSPID}
Dong, Z., Fan, X. \& Peng, Z., 2024. FNSPID: A Comprehensive Financial News
Dataset in Time Series. {\em Proceedings of the 30th ACM SIGKDD Conference on
Knowledge Discovery and Data Mining (KDD '24)}. arXiv:2402.06698.

\bibitem[Efron(1979)]{Efron1979Bootstrap}
Efron, B., 1979. Bootstrap Methods: Another Look at the Jackknife. {\em The
Annals of Statistics}, 7(1), pp. 1--26.

\bibitem[Gastinger et al.(2024)]{Gastinger2024Recurrence}
Gastinger, J. et al., 2024. History Repeats Itself: A Baseline for Temporal
Knowledge Graph Forecasting. {\em Proceedings of the 33rd International Joint
Conference on Artificial Intelligence (IJCAI)}. arXiv:2404.16726.

\bibitem[He et al.(2025)]{He2025Chrono}
He, S., Lv, L., Manela, A. \& Wu, J., 2025. Chronologically Consistent Large
Language Models. {\em arXiv preprint} arXiv:2502.21206.

\bibitem[Hu et al.(2022)]{Hu2022LoRA}
Hu, E. J. et al., 2022. LoRA: Low-Rank Adaptation of Large Language Models. {\em
International Conference on Learning Representations (ICLR)}. arXiv:2106.09685.

\bibitem[Jin et al.(2020)]{Jin2020RENET}
Jin, W., Qu, M., Jin, X. \& Ren, X., 2020. Recurrent Event Network:
Autoregressive Structure Inference over Temporal Knowledge Graphs. {\em
Proceedings of EMNLP}. arXiv:1904.05530.

\bibitem[Kleinberg(2003)]{Kleinberg2003Bursty}
Kleinberg, J., 2003. Bursty and Hierarchical Structure in Streams. {\em Data
Mining and Knowledge Discovery}, 7(4), pp. 373--397.

\bibitem[Kwon et al.(2023)]{Kwon2023vLLM}
Kwon, W. et al., 2023. Efficient Memory Management for Large Language Model
Serving with PagedAttention. {\em Proceedings of the 29th ACM Symposium on
Operating Systems Principles (SOSP)}. arXiv:2309.06180.

\bibitem[Li \& Sanna Passino(2024)]{Li2024FinDKG}
Li, X. V. \& Sanna Passino, F., 2024. FinDKG: Dynamic Knowledge Graphs with
Large Language Models for Detecting Global Trends in Financial Markets. {\em
Proceedings of the 5th ACM International Conference on AI in Finance (ICAIF
'24)}. arXiv:2407.10909.

\bibitem[Li et al.(2021)]{Li2021REGCN}
Li, Z. et al., 2021. Temporal Knowledge Graph Reasoning Based on Evolutional
Representation Learning. {\em Proceedings of the 44th International ACM SIGIR
Conference on Research and Development in Information Retrieval}. arXiv:2104.10353.

\bibitem[Liben-Nowell \& Kleinberg(2007)]{LibenNowell2007}
Liben-Nowell, D. \& Kleinberg, J., 2007. The Link-Prediction Problem for Social
Networks. {\em Journal of the American Society for Information Science and
Technology}, 58(7), pp. 1019--1031.

\bibitem[Page(1954)]{Page1954CUSUM}
Page, E. S., 1954. Continuous Inspection Schemes. {\em Biometrika}, 41(1/2),
pp. 100--115.

\bibitem[Qwen Team(2024)]{Qwen2024}
Qwen Team, 2024. Qwen2.5 Technical Report. {\em arXiv preprint} arXiv:2412.15115.

\bibitem[Trouillon et al.(2016)]{Trouillon2016ComplEx}
Trouillon, T. et al., 2016. Complex Embeddings for Simple Link Prediction. {\em
Proceedings of the 33rd International Conference on Machine Learning (ICML)}.
arXiv:1606.06357.

\bibitem[Vaswani et al.(2017)]{Vaswani2017Attention}
Vaswani, A. et al., 2017. Attention Is All You Need. {\em Advances in Neural
Information Processing Systems 30 (NeurIPS 2017)}. arXiv:1706.03762.

\bibitem[Vitter(1985)]{Vitter1985Reservoir}
Vitter, J. S., 1985. Random Sampling with a Reservoir. {\em ACM Transactions on
Mathematical Software}, 11(1), pp. 37--57.

\bibitem[Yang et al.(2015)]{Yang2015DistMult}
Yang, B. et al., 2015. Embedding Entities and Relations for Learning and
Inference in Knowledge Bases. {\em International Conference on Learning
Representations (ICLR)}. arXiv:1412.6575.

\end{thebibliography}

% =====================================================================
% APPENDICES
% =====================================================================
\appendix
\chapter{Reproducibility}
All experiments are scripted and seeded. Extraction ran with vLLM on two RTX
Pro 6000 GPUs (Chapter~\ref{ch:exp1}); every analysis, including the temporal
network with validation-based early stopping, the ChronoBERT scorers, the
bootstrap, the three burst detectors and the full theme system, reproduces on
a single 6\,GB consumer GPU. The analysis environment is Python 3.12 with
PyTorch (CUDA), networkx, NumPy, SciPy and pandas. Each graph construction is
a directory of FinDKG-format flat files (entity, relation and time indices;
chronologically split quadruple lists; weighted edge-week table); an
environment variable selects the construction so that every module runs
unchanged on any of the three. The derived knowledge-graph dataset (triples
and indices) and the source code are released for replication. Article text
is not re-redistributed: FNSPID is distributed under CC BY-NC 4.0 and the
underlying articles remain publisher-copyrighted, so the release instead
includes the exact FNSPID row identifiers of every sampled article together
with the seeded sampling scripts, from which both corpora reconstruct
exactly from the public FNSPID distribution. Fixed constants: weekly bucketing over 365 weeks; splits
0--254/255--309/310--364; detector gate $z \geq 3$, $O_t \geq 6$; windows
$r=4$, $b=26$ (defaults); Louvain seed 42; lifeline matching Jaccard
$\geq 0.3$ with an eight-week gap; NB dispersion $\alpha = 3.75$; bootstrap
$B = 1{,}000$; placebo draws $100{,}000$.
\chapter{Additional Results}
\section*{Cross-construction replication}
The method ranking of Chapter~\ref{ch:exp3} is identical on all three graph
constructions; absolute values are not comparable across candidate-set sizes.

\begin{table}[ht]
\centering
\caption{Best system and novelty across constructions.}
\begin{tabular}{lrrr}
\toprule
 & 200k core & 600k core & Dedup core \\
\midrule
Test novelty & 72.4\% & 64.0\% & 61.6\% \\
Recurrence MRR & 0.135 & 0.167 & 0.196 \\
Best backoff MRR / Hits@10 & 0.162 / 0.258 & 0.192 / 0.312 & 0.218 / 0.342 \\
Novel-edge ranking & \multicolumn{3}{c}{ChronoBERT $>$ RE-GCN $\geq$ ComplEx
$\gg$ CN/AA (all three)} \\
\bottomrule
\end{tabular}
\end{table}

\section*{Window-sensitivity grid}
Gold-list detections and placebo $p$-values for all nine window
configurations (recent $r$, baseline $b$, weeks):

\begin{table}[ht]
\centering
\caption{Theme detection across the $3\times3$ window grid.}
\begin{tabular}{lccc}
\toprule
 & $b=13$ & $b=26$ & $b=52$ \\
\midrule
$r=2$ & 7/10 ($p=3.9\times10^{-4}$) & 7/10 ($p=2.5\times10^{-4}$) &
7/10 ($p=2.9\times10^{-4}$) \\
$r=4$ & 8/10 ($p=1.2\times10^{-4}$) & 8/10 ($p=1.0\times10^{-4}$) &
8/10 ($p=3.2\times10^{-4}$) \\
$r=8$ & 8/10 ($p=3.8\times10^{-4}$) & 8/10 ($p=4\times10^{-5}$) &
8/10 ($p=4\times10^{-5}$) \\
\bottomrule
\end{tabular}
\end{table}

The $r=8$, $b=26$ cell achieves the earliest detection (median lead $-2$
weeks; 88\% early-or-on-time); shorter recent windows detect slightly later.

\end{document}